% GENERATED FILE. Edit canonical lessons, then run npm run generate:books.
% canonical-source-hash: fnv1a64:8763761b1a7fcdf5
% canonical-lessons: IT-C474-umore, IT-C474-sereno, IT-C474-lunatico, IT-C474-innervosirsi, IT-C474-sbuffare

\chapter{Di che umore sei?}
\label{ch:it-di-che-umore-sei}
\hlchaptermodality{474}

\begin{chapteropening}
\textbf{By the end of this chapter:} \emph{I can describe moods: calm, moody or on edge.}

Complete the last lesson of chapter 474: I can describe moods: calm, moody or on edge.
\end{chapteropening}

\section[l'umore]{l'umore — mood}

\label{lesson:IT-C474-umore}



\begin{quote}
Before the new one: say the Italian for sly, then the Italian for to admire.
\end{quote}

\subsection*{You'll want to know: l'umore}
\textbf{l'umore} — \textquotedblleft{}mood\textquotedblright{}.

\begin{quote}
Oggi sono di buon umore.
\end{quote}

\begin{quote}
\emph{Today I am in a good mood.}
\end{quote}

\emph{Di buon umore} — in a good mood; \emph{di cattivo umore} — in a bad mood.

\begin{cousinweb}
\emph{umore} is a learned borrowing from Latin \emph{ūmor} \textquotedblleft{}moisture, fluid\textquotedblright{} (EtymDB; Lewis \& Short): old medicine linked a person's mood to the fluids of the body.
\end{cousinweb}

\begin{grammarlens}[title={What you've built}]
One more word for telling, explaining and describing at B1.
\end{grammarlens}

\subsection*{Your turn}
\begin{itemize}
\raggedright
  \item \emph{Say it:} \emph{l'umore}
  \item \emph{Say it:} \emph{l'umore}, once more
  \item \emph{Say it:} a sentence of your own with \emph{l'umore}
  \item \emph{Recall:} say the Italian for sly, then the Italian for to admire, then say \emph{l'umore} again
\end{itemize}

\begin{tcolorbox}[breakable,colback=teal!4,colframe=teal!35!black,title={Before you move on}]
What does l'umore mean? (\textquotedblleft{}Mood\textquotedblright{}.) Say it once more.
\end{tcolorbox}
\section[sereno]{sereno — calm, serene; clear (of the sky)}

\label{lesson:IT-C474-sereno}



\begin{quote}
Before the new one: say the Italian for to admire, then the Italian for mood.
\end{quote}

\subsection*{You'll want to know: sereno}
\textbf{sereno} — \textquotedblleft{}calm, serene; clear (of the sky)\textquotedblright{}.

\begin{quote}
Dopo le vacanze mi sento più sereno.
\end{quote}

\begin{quote}
\emph{After the holidays I feel calmer.}
\end{quote}

Of the sky: \emph{il cielo è sereno} — the sky is clear.

\begin{cousinweb}
\emph{sereno} is inherited from Latin \emph{serēnus} \textquotedblleft{}clear, bright, calm\textquotedblright{} (EtymDB; Lewis \& Short), said first of the sky.
\end{cousinweb}

\begin{grammarlens}[title={What you've built}]
One more word for telling, explaining and describing at B1.
\end{grammarlens}

\subsection*{Your turn}
\begin{itemize}
\raggedright
  \item \emph{Say it:} \emph{sereno}
  \item \emph{Say it:} \emph{sereno}, once more
  \item \emph{Say it:} a sentence of your own with \emph{sereno}
  \item \emph{Recall:} say the Italian for to admire, then the Italian for mood, then say \emph{sereno} again
\end{itemize}

\begin{tcolorbox}[breakable,colback=teal!4,colframe=teal!35!black,title={Before you move on}]
What does sereno mean? (\textquotedblleft{}Calm, serene; clear (of the sky)\textquotedblright{}.) Say it once more.
\end{tcolorbox}
\section[lunatico]{lunatico — moody}

\label{lesson:IT-C474-lunatico}



\begin{quote}
Before the new one: say the Italian for mood, then the Italian for calm.
\end{quote}

\subsection*{You'll want to know: lunatico}
\textbf{lunatico} — \textquotedblleft{}moody\textquotedblright{}.

\begin{quote}
Mio fratello è lunatico: un giorno ride, un giorno no.
\end{quote}

\begin{quote}
\emph{My brother is moody: one day he laughs, the next he does not.}
\end{quote}

Not like English \emph{lunatic}: it means moody, full of ups and downs.

\begin{cousinweb}
\emph{lunatico} is a learned borrowing from Late Latin \emph{lūnāticus} \textquotedblleft{}moonstruck\textquotedblright{} (EtymDB), built on \emph{lūna} \textquotedblleft{}the moon\textquotedblright{}: moods were thought to follow its phases.
\end{cousinweb}

\begin{grammarlens}[title={What you've built}]
One more word for telling, explaining and describing at B1.
\end{grammarlens}

\subsection*{Your turn}
\begin{itemize}
\raggedright
  \item \emph{Say it:} \emph{lunatico}
  \item \emph{Say it:} \emph{lunatico}, once more
  \item \emph{Say it:} a sentence of your own with \emph{lunatico}
  \item \emph{Recall:} say the Italian for mood, then the Italian for calm, then say \emph{lunatico} again
\end{itemize}

\begin{tcolorbox}[breakable,colback=teal!4,colframe=teal!35!black,title={Before you move on}]
What does lunatico mean? (\textquotedblleft{}Moody\textquotedblright{}.) Say it once more.
\end{tcolorbox}
\section[innervosirsi]{innervosirsi — to get irritated, to get nervous}

\label{lesson:IT-C474-innervosirsi}



\begin{quote}
Before the new one: say the Italian for calm, then the Italian for moody.
\end{quote}

\subsection*{You'll want to know: innervosirsi}
\textbf{innervosirsi} — \textquotedblleft{}to get irritated, to get nervous\textquotedblright{}.

\begin{quote}
Mi innervosisco quando il treno è in ritardo.
\end{quote}

\begin{quote}
\emph{I get irritated when the train is late.}
\end{quote}

Built on \emph{nervoso}. An \emph{-isc-} verb with \emph{essere}: \emph{mi sono innervosito}.

\begin{cousinweb}
\emph{innervosirsi} is built on \emph{nervoso} \textquotedblleft{}nervous\textquotedblright{}, from Latin \emph{nervōsus} \textquotedblleft{}sinewy, vigorous\textquotedblright{} (Lewis \& Short), from \emph{nervus} \textquotedblleft{}a sinew\textquotedblright{}.
\end{cousinweb}

\begin{grammarlens}[title={What you've built}]
One more word for telling, explaining and describing at B1.
\end{grammarlens}

\subsection*{Your turn}
\begin{itemize}
\raggedright
  \item \emph{Say it:} \emph{innervosirsi}
  \item \emph{Say it:} \emph{innervosirsi}, once more
  \item \emph{Say it:} a sentence of your own with \emph{innervosirsi}
  \item \emph{Recall:} say the Italian for calm, then the Italian for moody, then say \emph{innervosirsi} again
\end{itemize}

\begin{tcolorbox}[breakable,colback=teal!4,colframe=teal!35!black,title={Before you move on}]
What does innervosirsi mean? (\textquotedblleft{}To get irritated, to get nervous\textquotedblright{}.) Say it once more.
\end{tcolorbox}
\section[sbuffare]{sbuffare — to huff, to puff with impatience}

\label{lesson:IT-C474-sbuffare}



\begin{quote}
Before the new one: say the Italian for moody, then the Italian for to get irritated.
\end{quote}

\subsection*{You'll want to know: sbuffare}
\textbf{sbuffare} — \textquotedblleft{}to huff, to puff with impatience\textquotedblright{}.

\begin{quote}
Il ragazzo sbuffa ogni volta che deve fare i compiti.
\end{quote}

\begin{quote}
\emph{The boy huffs every time he has to do his homework.}
\end{quote}

A puff of air through the lips; \emph{lo sbuffo} is the puff.

\begin{cousinweb}
\emph{sbuffare} is built from \emph{s-} and \emph{buffo} \textquotedblleft{}a puff of air\textquotedblright{}. The origin of \emph{buffo} is not settled; it is usually taken to imitate the sound of blowing.
\end{cousinweb}

\begin{grammarlens}[title={What you've built}]
One more word for telling, explaining and describing at B1.
\end{grammarlens}

\subsection*{Your turn}
\begin{itemize}
\raggedright
  \item \emph{Say it:} \emph{sbuffare}
  \item \emph{Say it:} \emph{sbuffare}, once more
  \item \emph{Say it:} a sentence of your own with \emph{sbuffare}
  \item \emph{Recall:} say the Italian for moody, then the Italian for to get irritated, then say \emph{sbuffare} again
\end{itemize}

\begin{tcolorbox}[breakable,colback=teal!4,colframe=teal!35!black,title={Before you move on}]
What does sbuffare mean? (\textquotedblleft{}To huff, to puff with impatience\textquotedblright{}.) Say it once more.
\end{tcolorbox}
